\begin{figure}[ht]
    \centering
    \hspace*{-2.0cm}
    \begin{tikzpicture}[
        ahlevel/.style={
                rectangle,
                rounded corners=5pt,
                minimum width=8cm,
                minimum height=1.2cm,
                text width=7cm,
                align=center,
                font=\sffamily\bfseries,
                draw=#1!60!black,
                top color=#1!15,
                bottom color=#1!35,
                text=#1!15!black,
                line width=0.7pt,
            },
        linklabel/.style={
                font=\sffamily\footnotesize\itshape,
                text=gray!50!black,
                fill=white,
                inner sep=2pt,
                rounded corners=2pt,
            },
        sidelabel/.style={
                font=\sffamily\footnotesize,
                text=#1,
                rotate=90,
            },
        desc/.style={
                font=\sffamily\footnotesize,
                text=gray!50!black,
                anchor=west,
                text width=4.0cm,
            },
        desc2/.style={
                font=\sffamily\normalsize,
                text=#1,
                anchor=west,
                text width=4.5cm,
            },
        ]

        % Hierarchy levels (top to bottom)
        \node[ahlevel=Orange](L5) at (-1.0, 5.4) {Funkčný účel};
        \node[ahlevel=Orange](L4) at (-1.0, 3.6) {Abstraktná funkcia};
        \node[ahlevel=Orange](L3) at (-1.0, 1.8) {Zovšeobecnená funkcia};
        \node[ahlevel=Green] (L2) at (-1.0, 0.0) {Fyzická funkcia};
        \node[ahlevel=Green] (L1) at (-1.0,-1.8) {Fyzická forma};

        % Descriptions (right side)
        \node[desc] at (3.0, 5.4)
        {Účel alebo ciele systému};
        \node[desc] at (3.0, 3.6)
        {Princípy toku, rovnováhy a riadenia};
        \node[desc] at (3.0, 1.8)
        {Procesy, interakcia komponentov};
        \node[desc] at (3.0, 0.0)
        {Schopnosti, zaradenia a komponenty};
        \node[desc] at (3.0,-1.8)
        {Dimenzie a fyzikálne vlastnosti prostredia};

        \draw[loosely dotted, line width=2pt, color=black]
        (-6.0, 0.9) -- (11.6, 0.9);


        % left Side labels
        \node[sidelabel=Orange, anchor=south] at (-5.2, 3.6)  {koncepčná rovina};
        \node[sidelabel=Green, anchor=south] at (-5.2,-0.9)  {skutočná/reálna rovina};

        % left Side arrows (means-ends)
        \draw[{Stealth[length=2.5mm]}-{Stealth[length=2.5mm]}, line width=0.8pt, color=Orange!40]
        (-5.2, 6.3) -- (-5.2, 0.9);
        \draw[{Stealth[length=2.5mm]}-{Stealth[length=2.5mm]}, line width=0.8pt, color=Green!40]
        (-5.2, 0.9) -- (-5.2, -2.7);

        % right Side arrows (means-ends)
        \draw[{Stealth[length=2.5mm]}-{Stealth[length=2.5mm]}, line width=0.8pt, color=Orange!40]
        (7.5, 6.3) -- (7.5, 0.9);
        \draw[{Stealth[length=2.5mm]}-{Stealth[length=2.5mm]}, line width=0.8pt, color=Green!40]
        (7.5, 0.9) -- (7.5, -2.7);

        % left side labels
        \node[desc2=Orange] at (7.5, 3.6)
        {\textbf{systemovo-orientovaná analýza rizík}};
        \node[desc2=Green] at (7.5, -0.9)
        {\textbf{komponentovo-orientovaná analýza rizík}};

    \end{tikzpicture}
    \caption{Rasmussenova hierarchia abstrakcie}
    \label{fig:abstraction-hierarchy}
\end{figure}